\documentclass[10pt, a4paper]{article}

%% ========== Page geometry ==========
\usepackage[
    a4paper,
    ignoreheadfoot,
    top=1.5cm, bottom=1.5cm,
    left=1.8cm, right=1.8cm,
    footskip=0.8cm
]{geometry}

%% ========== XeLaTeX fonts ==========
\usepackage{fontspec}
\usepackage{xeCJK}
\defaultfontfeatures{Ligatures=TeX}
\setmainfont{Helvetica Neue}
\setsansfont{Helvetica Neue}
\setCJKmainfont{PingFang SC}
\setCJKsansfont{PingFang SC}

%% ========== Core packages ==========
\usepackage[explicit]{titlesec}
\usepackage{tabularx}
\usepackage{array}
\usepackage[dvipsnames]{xcolor}
\usepackage{enumitem}
\usepackage{fontawesome5}
\usepackage{graphicx}
\usepackage{tikz}
\usepackage[pscoord]{eso-pic}
\usepackage{calc}
\usepackage{lastpage}
\usepackage{changepage}
\usepackage{paracol}
\usepackage{ifthen}
\usepackage{needspace}
\usepackage{etoolbox}

%% ========== Color palette ==========
\definecolor{primaryColor}{RGB}{0, 79, 144}
\definecolor{accentBg}{RGB}{232, 240, 254}
\definecolor{textColor}{RGB}{45, 45, 45}
\definecolor{grayText}{RGB}{107, 114, 128}
\definecolor{ruleColor}{RGB}{209, 213, 219}

%% ========== Hyperlinks (must precede bookmark) ==========
\usepackage[
    pdftitle={Haohan Lin's CV},
    pdfauthor={Haohan Lin},
    colorlinks=true,
    urlcolor=primaryColor,
    linkcolor=primaryColor
]{hyperref}
\usepackage{bookmark}

%% ========== Global settings ==========
\color{textColor}
\AtBeginEnvironment{adjustwidth}{\partopsep0pt}
\pagestyle{empty}
\setcounter{secnumdepth}{0}
\setlength{\parindent}{0pt}
\setlength{\topskip}{0pt}
\setlength{\columnsep}{0.15cm}
\linespread{1.05}

%% ========== Footer ==========
\makeatletter
\let\ps@cvFooter\ps@plain
\patchcmd{\ps@cvFooter}{\thepage}{%
    \color{grayText}\textit{\small Haohan Lin — Page \thepage{} of \pageref*{LastPage}}%
}{}{}
\makeatother
\pagestyle{cvFooter}

%% ========== Top accent bar (every page) ==========
\AddToShipoutPictureBG{%
    \AtPageUpperLeft{%
        \raisebox{-3pt}[0pt][0pt]{\color{primaryColor}\rule{\paperwidth}{3pt}}%
    }%
}

%% ========== Section formatting ==========
\titleformat{\section}{
    \needspace{4\baselineskip}
    \Large\color{primaryColor}
}{}{0pt}{
    \textbf{#1}\hspace{0.15cm}{\color{primaryColor}\titlerule[0.6pt]}
}[]

\titlespacing{\section}{-1pt}{0.35cm}{0.2cm}

%% ========== Custom environments ==========

\newenvironment{highlights}{
    \begin{itemize}[
        topsep=0.10cm, parsep=0.10cm,
        partopsep=0pt, itemsep=0pt,
        leftmargin=0.4cm + 10pt
    ]
}{
    \end{itemize}
}

\newenvironment{highlightsforbulletentries}{
    \begin{itemize}[
        topsep=0.10cm, parsep=0.10cm,
        partopsep=0pt, itemsep=0pt,
        leftmargin=10pt
    ]
}{
    \end{itemize}
}

\newenvironment{onecolentry}{
    \begin{adjustwidth}{0.2cm}{0.2cm}
}{
    \end{adjustwidth}
}

\newenvironment{twocolentry}[2][]{
    \onecolentry
    \def\secondColumn{#2}
    \setcolumnwidth{\fill, 4.5cm}
    \begin{paracol}{2}
}{
    \switchcolumn \raggedleft \secondColumn
    \end{paracol}
    \endonecolentry
}

\newenvironment{threecolentry}[3][]{
    \onecolentry
    \def\thirdColumn{#3}
    \setcolumnwidth{1cm, \fill, 4.5cm}
    \begin{paracol}{3}
    {\raggedright #2} \switchcolumn
}{
    \switchcolumn \raggedleft \thirdColumn
    \end{paracol}
    \endonecolentry
}

%% ========== Photo commands ==========
\newcommand{\cvPhotoSize}{2.55cm}
% Square photo (no frame): \photo[size]{filename}
\newcommand{\photo}[2][2.9cm]{%
    \begin{tikzpicture}
        \pgfmathsetlengthmacro\photohalf{#1/2}%
        \begin{scope}
            \clip (-\photohalf,-\photohalf) rectangle (\photohalf,\photohalf);
            \node[inner sep=0pt] at (0,0) {\includegraphics[width=#1,height=#1,keepaspectratio]{#2}};
        \end{scope}
    \end{tikzpicture}%
}

% Placeholder when no photo file exists (no frame)
\newcommand{\photoplaceholder}[1][2.9cm]{%
    \begin{tikzpicture}
        \node[inner sep=0pt, text=grayText, font=\Large] at (0,0.12) {\faUser};
        \node[inner sep=0pt, text=grayText, font=\tiny] at (0,-0.45) {Photo};
    \end{tikzpicture}%
}

%% ========== Link with external-link icon ==========
\let\hrefWithoutArrow\href
\renewcommand{\href}[2]{%
    \hrefWithoutArrow{#1}{%
        \ifthenelse{\equal{#2}{}}{ }{#2 }\raisebox{.15ex}{\footnotesize \faExternalLink*}%
    }%
}

%% ========== Contact separator ==========
\newcommand{\contactsep}{\kern 0.25cm{\color{ruleColor}\textbar}\kern 0.25cm}

%% ========== Last-updated overlay ==========
\newcommand{\placelastupdatedtext}{%
    \AddToShipoutPictureFG*{%
        \put(
            \LenToUnit{\paperwidth-2cm-0.2cm+0.05cm},
            \LenToUnit{\paperheight-1.0cm}
        ){\vtop{{\null}\makebox[0pt][c]{
            \small\color{grayText}\textit{Last updated in April 2026}\hspace{\widthof{Last updated in April 2026}}
        }}}%
    }%
}

% ================================================================
%  DOCUMENT BEGIN
% ================================================================
\begin{document}

%% ==========================================================
%%  中文简历
%% ==========================================================

% --- 中文页头（含照片） ---
\begin{minipage}[t]{\dimexpr\textwidth-3.4cm\relax}
    % >>> 如需显示中文名，将下行改为: \textbf{林浩翰 \;|\; Haohan Lin}
    {\fontsize{28pt}{28pt}\selectfont\color{primaryColor}\textbf{林浩瀚}\par}

    % \vspace{0.12cm}

    % {\small\color{grayText}软件开发 \textperiodcentered\ 自动化测试 \textperiodcentered\ Linux / Shell 运维方向\par}

    \vspace{0.18cm}

    {\small\color{textColor}%
        \mbox{\faMapMarker*\hspace{0.1cm}苏州}%
        \contactsep
        \mbox{\hrefWithoutArrow{mailto:guameezia@gmail.com}{\faEnvelope[regular]\hspace{0.1cm}guameezia@gmail.com}}%
        \contactsep
        \mbox{\hrefWithoutArrow{tel:+86 130 2388 0536}{\faPhone*\hspace{0.1cm}130 2388 0536}}%

        \vspace{0.12cm}

        \mbox{\hrefWithoutArrow{https://www.linkedin.com/in/haohan-lin-a861aa359/}{\faLinkedinIn\hspace{0.1cm}Haohan Lin}}%
        \contactsep
        \mbox{\hrefWithoutArrow{https://github.com/Guameezia}{\faGithub\hspace{0.1cm}Guameezia}}%
    }
\end{minipage}%
\hfill
\begin{minipage}[t]{3.0cm}
    \vspace{0pt}\raggedleft
    \photo[\cvPhotoSize]{profile.png}
\end{minipage}

\vspace{-0.14cm}
\noindent{\color{ruleColor}\hrule height 0.4pt}
\vspace{0.12cm}

\vspace{-0.3cm}
% --- 教育经历 ---
\section{教育经历}

    \begin{threecolentry}{\textbf{本科}}{
        中国苏州

        2024.9 – 2028.7
    }
        \textbf{信息与计算科学专业（计算机科学）} \\ 西交利物浦大学（XJTLU）
        \begin{highlights}
            \item \textbf{核心课程：} 数据结构、操作系统、计算机系统、计算机网络、数据库、算法、软件工程、Java 程序设计、人工智能、机器学习、人工智能应用设计、多智能体系统、人机交互、计算机图形学、离散数学与统计学、线性代数、多元微积分
            \item \textbf{第一学年平均成绩：} 74\%
            \item \textbf{第二学年平均成绩（第一学期）：} 78\%（专业前 20\%）
            \item[] {\small\color{grayText}\textit{英制百分制（0--100）；70 分及以上为一等区间（成绩上限档，相当于满绩）}}
        \end{highlights}
    \end{threecolentry}

\vspace{-0.3cm}

% --- 项目经历 ---
\section{项目经历}

    \begin{twocolentry}{
        中国苏州

        2025.6 – 2025.8

        \vspace{1em}

        \href{https://github.com/Guameezia/2025_SURF_Poster}{github.com/Guameezia/ \\ 2025\_SURF\_Poster}
    }
        \textbf{SVD 与 CNN 在图像处理中的对比与优化}
        \begin{highlights}
            \item 导师：杨其儒博士；基于 Python（NumPy / PyTorch / Matplotlib）搭建图像处理与机器学习实验流水线，覆盖数据预处理、训练与评估全流程
            \item 设计 SVD 低秩近似与 CNN 两套图像重建方案的对比实验 %，从压缩比、重建误差（PSNR）与运行效率多维度量化分析，输出可复现的实验脚本与结果
            \item 使用 Git 进行版本管理与多人协作，任科研海报 LaTeX 总编辑%，负责终版版式统一；项目获学校正式科研证明
        \end{highlights}
    \end{twocolentry}

    \vspace{0.2cm}

    \begin{twocolentry}{
        中国苏州

    2026.3 – 至今
    }
        \textbf{AI Agent 驱动的校园室内导航与实时社交一体化平台（UniNavi）}
        \begin{highlights}
            \item 导师：马腾博士（在研）；负责室内导航与 AI Agent 模块的设计与工程实现
            \item 搭建多层室内路径规划，开发具身 AI Agent 数字人，基于 LLM 与 RAG 检索增强构建对话引擎，结合 Function Calling 工具调用与 ReAct 推理规划完成校园导引与多任务编排。Git 管理版本迭代，持续完善功能与稳定性
            % \item 采用前后端分离架构（React 前端 + Python 服务端），通过 \textbf{Docker} 容器化部署到 \textbf{Linux} 服务器，统一开发与运行环境
            % \item 使用 Git 管理多人协作与版本迭代，持续完善功能与稳定性
        \end{highlights}
    \end{twocolentry}

\vspace{-0.3cm}

% --- 工作经历 ---
\section{工作经历}

    \begin{twocolentry}{
        中国上海

        2026.1 – 2026.3
    }
        \textbf{测试开发实习生}，上海兑观信息科技有限公司
        \begin{highlights}
            \item 基于 Shell/Bash 编写自动化测试与运维辅助脚本%（数据提取、日志清洗、批量校验），并用 \textbf{cron} 配置定时任务，将多项重复性操作脚本化，单次执行时间从约 30 分钟降至 5 分钟内
            % \item 在 Linux 服务器上排查服务异常：查看系统与应用日志、定位异常进程与端口占用、复现并回归验证缺陷，协助保障测试环境与系统稳定运行
            \item 使用 Docker 拉起与复现测试环境，统一本地与服务器依赖版本%，缩短环境搭建与问题复现时间
            % \item 执行功能、性能与接口测试，累计编写并维护 60+ 条测试用例，提交与跟踪 40+ 个缺陷，输出可复现的缺陷报告与步骤
            \item 使用 Python（Playwright） 编写爬虫脚本批量抓取分析平台数据，清洗后结合 Excel 元数据对接公司系统数据库，支撑测试数据准备与功能验证
        \end{highlights}
    \end{twocolentry}

\vspace{-0.3cm}

% --- 技能特长 ---
\section{技能特长}

    \begin{onecolentry}
        \textbf{编程与标记语言：} C/C++、Python、Java、JavaScript、TypeScript、SQL、HTML、CSS、MATLAB、LaTeX
    \end{onecolentry}

    \vspace{0.1cm}

    \begin{onecolentry}
        \textbf{运维与系统：} Linux（Ubuntu）、Shell / Bash 脚本、日志排查与进程/端口定位、Docker 容器化、Git 版本管理
    \end{onecolentry}

    \vspace{0.1cm}

    \begin{onecolentry}
        \textbf{测试与质量：} 功能 / 性能 / 接口测试、Playwright 自动化脚本、缺陷管理与回归验证、测试用例设计
    \end{onecolentry}

    \vspace{0.1cm}

    \begin{onecolentry}
        \textbf{工具与平台：} Cursor / VS Code / Xcode 等 IDE，熟练运用 AI 辅助编程（Cursor、Codex、Claude Code）提升开发效率；React 前端开发；macOS / Linux 双平台；Microsoft Office 基础操作
    \end{onecolentry}

% --- 个人总结 ---
\section{个人总结}

    \begin{onecolentry}
        \textbf{计算机科学}方向本科在读（西交利物浦大学，信息与计算科学），具备扎实的编程基础与工程实践能力。熟悉 \textbf{Linux} 环境下的 \textbf{Shell/Bash} 脚本编写、日志排查与 \textbf{Docker} 容器化部署，并有自动化测试与质量保障经验；兼具 SURF 科研、UniNavi 在研项目与测试开发实习等经历，能快速上手新技术。
    \end{onecolentry}

\clearpage


%% ==========================================================
%%  ENGLISH CV
%% ==========================================================
\begingroup
\small
\linespread{0.96}\selectfont
\emergencystretch=1em
\titlespacing{\section}{-1pt}{0.22cm}{0.10cm}
\renewenvironment{highlights}{
    \begin{itemize}[
        topsep=0.04cm, parsep=0.02cm,
        partopsep=0pt, itemsep=0pt,
        leftmargin=0.34cm + 8pt
    ]
}{
    \end{itemize}
}
\renewenvironment{onecolentry}{
    \begin{adjustwidth}{0.12cm}{0.12cm}
}{
    \end{adjustwidth}
}
\renewenvironment{twocolentry}[2][]{
    \onecolentry
    \def\secondColumn{#2}
    \setcolumnwidth{\fill, 3.8cm}
    \begin{paracol}{2}
}{
    \switchcolumn \raggedleft \secondColumn
    \end{paracol}
    \endonecolentry
}
\renewenvironment{threecolentry}[3][]{
    \onecolentry
    \def\thirdColumn{#3}
    \setcolumnwidth{0.9cm, \fill, 3.8cm}
    \begin{paracol}{3}
    {\raggedright #2} \switchcolumn
}{
    \switchcolumn \raggedleft \thirdColumn
    \end{paracol}
    \endonecolentry
}

% --- Header with photo (top-aligned; gray rule sits slightly up toward photo) ---
\begin{minipage}[t]{\dimexpr\textwidth-3.4cm\relax}
    {\fontsize{25pt}{25pt}\selectfont\color{primaryColor}\textbf{Haohan Lin}\par}

    % \vspace{0.12cm}

    % {\small\color{grayText}Software Development \textperiodcentered\ Test Automation \textperiodcentered\ Linux / Shell Operations\par}

    \vspace{0.10cm}

    {\small\color{textColor}%
        \mbox{\faMapMarker*\hspace{0.1cm}Suzhou, China}%
        \contactsep
        \mbox{\hrefWithoutArrow{mailto:guameezia@gmail.com}{\faEnvelope[regular]\hspace{0.1cm}guameezia@gmail.com}}%
        \contactsep
        \mbox{\hrefWithoutArrow{tel:+86 130 2388 0536}{\faPhone*\hspace{0.1cm}+86 130 2388 0536}}%

        \vspace{0.06cm}

        \mbox{\hrefWithoutArrow{https://www.linkedin.com/in/haohan-lin-a861aa359/}{\faLinkedinIn\hspace{0.1cm}Haohan Lin}}%
        \contactsep
        \mbox{\hrefWithoutArrow{https://github.com/Guameezia}{\faGithub\hspace{0.1cm}Guameezia}}%
    }
\end{minipage}%
\hfill
\begin{minipage}[t]{3.0cm}
    \vspace{0pt}\raggedleft
    \photo[\cvPhotoSize]{profile.png}
\end{minipage}

\vspace{-0.14cm}
\noindent{\color{ruleColor}\hrule height 0.4pt}
\vspace{0.06cm}

\vspace{-0.2cm}
% --- Education ---
\section{Education}

    \begin{threecolentry}{\textbf{BSc}}{
        Suzhou, China

        Sept 2024 – Jul 2028
    }
        \textbf{Information and Computing Science} \\ Xi'an Jiaotong-Liverpool University (XJTLU)
        \begin{highlights}
            \item \textbf{Core coursework:} OS, Computer Systems, Networks, Databases, Data Structures, Algorithms, Software Engineering, Java, AI, ML, Multi-Agent Systems, HCI, Graphics, Discrete Math/Stats, Linear Algebra, Multivariable Calculus
            \item \textbf{Stage 1 average:} 74\%
            \item \textbf{Stage 2 average (Semester 1):} 78\% (top 20\% in programme)
            \item[] {\small\color{grayText}\textit{UK-style marks (0--100); 70+ = First Class (top band; near full marks).}}
            % \item \textbf{Current main skill course grades:} \\
            % CPT001 Professional Skills in Computer Science \hfill 82\% \\
            % MTH007 Linear Algebra \hfill 94\% \\
            % MTH004 Multivariable Calculus \hfill 96\% \\
            % (Can provide complete transcript)
        \end{highlights}
    \end{threecolentry}

\vspace{-0.2cm}
% --- Projects ---
\section{Projects}

    \begin{twocolentry}{
        Suzhou, China

    Jun 2025 – Aug 2025

    \vspace{0.45em}

    \href{https://github.com/Guameezia/2025_SURF_Poster}{github.com/Guameezia/ \\ 2025\_SURF\_Poster}
    }
        \textbf{Comparative Study and Optimization of SVD and CNN in Image Processing}
        \begin{highlights}
            \item Supervised by Dr.\ Chi-Ru Yang; built an image-processing and ML pipeline in \textbf{Python} (NumPy / PyTorch / Matplotlib), covering pre-processing, training, and evaluation
            \item Designed comparison experiments between SVD low-rank approximation and CNN reconstruction %, quantifying compression ratio, reconstruction error (PSNR), and runtime efficiency with reproducible scripts and results
            \item Used Git for version control and team collaboration; served as chief LaTeX editor of the research poster %; responsible for the final layout; project received a university-issued research certificate
        \end{highlights}
    \end{twocolentry}

    \vspace{0.06cm}

    \begin{twocolentry}{
        Suzhou, China

    Mar 2026 – Present
    }
        \textbf{AI-Agent-Driven Integrated Campus Indoor Navigation and Real-Time Social Platform (UniNavi)}
        \begin{highlights}
            \item Supervised by Dr.\ Teng Ma (ongoing); responsible for indoor-navigation and AI-agent modules from design to implementation
            \item Building multi-floor indoor path planning and an embodied AI-agent avatar; developed an \textbf{LLM} + \textbf{RAG} dialogue engine with Function Calling and ReAct reasoning for campus guidance and multi-task orchestration; managing Git-based iteration
            % \item Front-end/back-end split architecture (React front end + Python back end), containerized with \textbf{Docker} and deployed to a \textbf{Linux} server for consistent dev and runtime environments
            % \item Managing team collaboration and version iteration with Git while continuously improving features and stability
        \end{highlights}
    \end{twocolentry}

% --- Work Experience ---
\vspace{-0.2cm}
\section{Work Experience}

    \begin{twocolentry}{
        Shanghai, China

    Jan 2026 – Mar 2026
    }
        \textbf{Test Development Intern}, Shanghai Videt Information Technology Co., Ltd.
        \begin{highlights}
            \item Wrote \textbf{Shell/Bash} scripts for automated testing and operations support %(data extraction, log cleaning, batch validation) and scheduled them with \textbf{cron}, cutting recurring manual tasks from $\sim$30 min to under 5 min per run
            % \item Troubleshot service issues on \textbf{Linux} servers: inspected system/application logs, traced abnormal processes and port usage, reproduced and regression-verified defects to help keep test environments and systems stable
            \item Used \textbf{Docker} to spin up and reproduce test environments, aligning dependency versions across local and server setups %and reducing setup/repro time
            % \item Ran functional, performance, and API testing; authored and maintained 60+ test cases and filed/tracked 40+ defects with reproducible bug reports
            \item Wrote \textbf{Python} (Playwright) crawler scripts to batch-scrape analysis-platform data, clean it, and integrate it with Excel metadata and the company's system database for test-data preparation and functional verification
        \end{highlights}
    \end{twocolentry}

% --- Skills ---
\section{Skills}

    \begin{onecolentry}
        \textbf{Programming \& Markup:} C/C++, Python, Java, JavaScript, TypeScript, SQL, HTML, CSS, MATLAB, LaTeX
    \end{onecolentry}

\vspace{0.06cm}

    \begin{onecolentry}
        \textbf{Operations \& Systems:} Linux (Ubuntu), Shell / Bash scripting, log troubleshooting and process/port diagnosis, Docker containerization, Git version control
    \end{onecolentry}

\vspace{0.06cm}

    \begin{onecolentry}
        \textbf{Testing \& Quality:} functional / performance / API testing, Playwright automation, defect management and regression, test-case design
    \end{onecolentry}

\vspace{0.06cm}

    \begin{onecolentry}
        \textbf{Tools \& Platforms:} IDEs such as Cursor / VS Code / Xcode, proficient with AI-assisted coding (Cursor, Codex, Claude Code) to boost development efficiency; React front-end development; macOS / Linux dual platform; basic Microsoft Office
    \end{onecolentry}

% --- Summary ---
\section{Summary}

    \begin{onecolentry}
        Undergraduate \textbf{computer science} student (Information and Computing Science, XJTLU) with a solid programming foundation and hands-on engineering experience. Comfortable with \textbf{Shell/Bash}, log troubleshooting, \textbf{Docker} on \textbf{Linux}, test automation, and quality assurance; experience spans SURF research, UniNavi, and a test-development internship.
    \end{onecolentry}

\endgroup
\end{document}
